\chapter{Signal Processing and Numerical Foundation}
\label{ch:signal-processing}

\section{Signal-Processing Gate Roadmap}
This chapter contains the complete SP-G0 through SP-G7 roadmap. The chapter is written
before all experiments are complete so that future results fill predefined evidence locations
rather than changing the scientific question after outcomes are known.

\section{SP-G0 --- Signal-Processing Contract Freeze}
\label{sec:sp-g0}

\subsection{Purpose}

SP-G0 establishes the authoritative signal-processing contract that all later
software, fixed-point, HLS, RTL, and board-level verification shall use.

The purpose of this gate is to prevent numerical or hardware implementation
from proceeding while basic array-processing conventions remain implicit or
ambiguous.

The central requirement is that the same physical and numerical interpretation
of the input snapshot matrix shall be preserved across the complete chain

\[
\text{Signal Model}
\rightarrow
\text{Floating-Point Reference}
\rightarrow
\text{Fixed-Point Reference}
\rightarrow
\text{HLS/RTL}
\rightarrow
\text{Board}.
\]

\subsection{Engineering Question}

The engineering question for SP-G0 is:

\begin{quote}
What exact signal, array, indexing, normalization, complex-data, and frame
conventions define the canonical input and covariance output of the project?
\end{quote}

These conventions must be frozen before SP-G1 begins.

\subsection{Currently Frozen Baseline}

The following baseline properties are already established by the approved
project scope and therefore do not require reinvention:

\begin{itemize}
    \item number of sensor channels:
    \[
    N = 8;
    \]

    \item default number of snapshots:
    \[
    M = 128;
    \]

    \item snapshot matrix:
    \[
    \mathbf X \in \mathbb C^{N\times M};
    \]

    \item snapshots are represented as columns of $\mathbf X$;

    \item the sample covariance matrix is defined as
    \[
    \widehat{\mathbf R}
    =
    \frac{1}{M}
    \mathbf X \mathbf X^{\mathsf H};
    \]

    \item the baseline FPGA responsibility is sample covariance matrix
    formation;

    \item downstream DOA estimation remains in software for the baseline
    project;

    \item the baseline geometry belongs to the uniform linear array family;

    \item controlled synthetic data are used for baseline verification;

    \item verification shall include at least one deterministic reference case
    together with a small controlled regression set;

    \item the existing floating-point implementation is treated as prior
    engineering evidence, but its conventions must still be audited before
    they become the formal signal-processing contract.
\end{itemize}

\subsection{Canonical ULA and Steering Contract}

The SP-G0 source audit showed that the predecessor HLS implementation
provides direct evidence for the SCM arithmetic and deterministic vector, but
does not independently define every physical array-manifold convention needed
by this project.

The remaining physical conventions are therefore frozen explicitly for this
repository.

The canonical baseline is an eight-element uniform linear array with
wavelength-normalized sensor coordinates

\[
\xi_n
=
\frac{x_n}{\lambda}
=
\frac{n}{2},
\qquad
n=0,\ldots,7.
\]

Thus,

\[
\frac{d}{\lambda}=0.5.
\]

The canonical array axis is \(+x\). The one-parameter visible sector is

\[
-90^\circ
\leq
\theta
\leq
+90^\circ,
\]

with

\[
\theta=0
\]

defined as broadside and positive \(\theta\) rotating toward \(+x\).

For the receive-manifold derivation, the project adopts the temporal phasor
convention

\[
\exp(-j\omega t).
\]

The canonical receive steering coefficient is

\[
a_n(\theta)
=
\exp\!\left[
-j\,2\pi\,\xi_n\sin\theta
\right].
\]

Equivalently,

\[
a_n(\theta)
=
\exp\!\left[
-jk(x_n-x_{\mathrm{ref}})\sin\theta
\right].
\]

For the baseline golden case,

\[
x_{\mathrm{ref}}=x_0=0.
\]

This first-sensor phase reference is a deliberate project decision and shall
not be silently replaced by another reference convention in downstream
software or hardware.

\subsection{Deterministic Golden Reference Case}

The canonical deterministic SP-G0 case is

\[
N=8,
\qquad
M=128,
\qquad
\theta=+30^\circ,
\]

with unit source amplitude, zero source phase, and no additive noise.

The source waveform is

\[
s[m]=1+j0,
\qquad
m=0,\ldots,127.
\]

Therefore,

\[
\mathbf X[:,m]
=
\mathbf a(30^\circ)
\]

for every snapshot.

Because

\[
\frac{d}{\lambda}=0.5
\]

and

\[
\sin 30^\circ=\frac{1}{2},
\]

the steering coefficient becomes

\[
a_n
=
\exp\!\left(-j\frac{\pi n}{2}\right),
\]

which gives

\[
\mathbf a
=
\begin{bmatrix}
1 &
-j &
-1 &
+j &
1 &
-j &
-1 &
+j
\end{bmatrix}^{\mathsf T}.
\]

This vector exactly reproduces the deterministic vector used by the
predecessor HLS testbench.

Because all \(M\) snapshots are identical,

\[
\mathbf X
=
\mathbf a\,\mathbf 1_M^{\mathsf T},
\]

and the resulting deterministic golden SCM is

\[
\widehat{\mathbf R}
=
\frac{1}{M}
\mathbf X\mathbf X^{\mathsf H}
=
\mathbf a\mathbf a^{\mathsf H}.
\]

Consequently,

\[
\widehat R_{ii}=1
\]

and

\[
\widehat{\mathbf R}
=
\widehat{\mathbf R}^{\mathsf H}.
\]

This analytic matrix becomes an independent structural reference for the
later floating-point golden implementation.

\subsection{Mean-Treatment Contract}

The predecessor HLS operation performs no sample-mean subtraction.

The canonical operation is therefore

\[
\widehat{\mathbf R}
=
\frac{1}{M}
\mathbf X\mathbf X^{\mathsf H}.
\]

In strict statistical terminology this is an uncentered sample second-moment
matrix when the sample mean is nonzero.

The project uses the conventional array-processing term sample covariance
matrix (SCM) for this operation, while explicitly preserving the fact that no
mean-centering is performed.

A mean-centered estimator shall not silently replace this baseline operation.

\subsection{Regression and Logical-Interface Contract}

Randomized regression shall use the explicit NumPy Generator API with

\[
\texttt{PCG64}
\]

and root seed

\[
0.
\]

Independent randomized cases shall use deterministic
\texttt{SeedSequence}-derived substreams rather than implicit global random
state.

The logical input sample is

\[
X[n,m],
\]

where \(n\) is the sensor/channel index and \(m\) is the snapshot index.

Thus,

\[
\mathbf X\in\mathbb C^{8\times128}
\]

with sensors as rows and snapshots as columns.

This definition freezes the signal-processing frame semantics only.

Physical AXI serialization, including real/imaginary packing, byte order,
word width, sign extension, stream traversal, DMA framing, and TLAST
behavior, belongs to INT-G0.

\subsection{Evidence Classification and Convention Conflict}

SP-G0 distinguishes four evidence classes:

\begin{itemize}
    \item \textbf{Source evidence}: behavior explicitly present in the
    predecessor controlled implementation;

    \item \textbf{Derived consistency check}: a mathematical consequence
    shown to be compatible with source evidence;

    \item \textbf{Cross-project precedent}: a convention established in
    another controlled project but not automatically authoritative here;

    \item \textbf{New project decision}: a convention explicitly selected
    where predecessor evidence was insufficient.
\end{itemize}

This distinction is necessary because historical software exists with both
positive- and negative-sign spatial exponential conventions.

For this repository, the authoritative receive steering convention is

\[
a_n(\theta)
=
\exp\!\left[
-j\,2\pi\,\xi_n\sin\theta
\right].
\]

Historical positive-sign implementations are not declared universally
incorrect in isolation, because steering sign depends on the complete
temporal-phasor, direction, and receive-model convention.

They are, however, non-canonical for this repository unless an explicit
convention-adaptation and verification step is performed.
\subsection{Covariance Definition}

For the approved frame

\[
\mathbf X
=
\begin{bmatrix}
x_1[0] & x_1[1] & \cdots & x_1[M-1] \\
x_2[0] & x_2[1] & \cdots & x_2[M-1] \\
\vdots & \vdots & \ddots & \vdots \\
x_N[0] & x_N[1] & \cdots & x_N[M-1]
\end{bmatrix},
\]

the covariance element is

\[
\widehat R_{ij}
=
\frac{1}{M}
\sum_{m=0}^{M-1}
x_i[m]\,x_j^{*}[m].
\]

Therefore,

\[
\widehat{\mathbf R}
=
\frac{1}{M}
\mathbf X\mathbf X^{\mathsf H}.
\]

The conjugation convention is fundamental because the FPGA implementation
must reproduce the same complex multiplication rule as the numerical
reference.

For

\[
x_i = a + jb,
\qquad
x_j = c + jd,
\]

the required product is

\[
x_i x_j^{*}
=
(a+jb)(c-jd)
=
(ac+bd)
+
j(bc-ad).
\]

This expression shall later be used directly in fixed-point range analysis
and HLS verification.

\subsection{Statistical and Numerical Boundaries}

SP-G0 also establishes the terminology required to avoid mixing distinct
sources of error.

The project shall preserve the hierarchy

\[
\mathbf R_x
\rightarrow
\widehat{\mathbf R}_{fp}
\rightarrow
\widehat{\mathbf R}_{fxp}
\rightarrow
\widehat{\mathbf R}_{hw},
\]

where

\begin{itemize}
    \item $\mathbf R_x$ denotes the population covariance model where such a
    model is defined;

    \item $\widehat{\mathbf R}_{fp}$ denotes the finite-snapshot
    floating-point sample covariance matrix;

    \item $\widehat{\mathbf R}_{fxp}$ denotes the bit-accurate fixed-point
    numerical result;

    \item $\widehat{\mathbf R}_{hw}$ denotes a hardware-produced result,
    including HLS/RTL or board output as explicitly identified.
\end{itemize}

Consequently,

\[
\widehat{\mathbf R}_{fp} - \mathbf R_x
\]

is associated with finite-snapshot statistical estimation behavior, while

\[
\widehat{\mathbf R}_{fxp}
-
\widehat{\mathbf R}_{fp}
\]

is associated with finite-word-length numerical behavior.

These quantities shall not be presented as the same error source.

\subsection{SP-G0 Evidence Closure}

SP-G0 produced traceable evidence for:

\begin{enumerate}
    \item predecessor HLS matrix dimensions and indexing;

    \item covariance normalization and conjugation behavior;

    \item absence of sample-mean subtraction;

    \item the predecessor deterministic steering vector;

    \item mathematical consistency of that vector with the adopted
    half-wavelength ULA and steering convention;

    \item explicit resolution of the historical steering-sign conflict;

    \item the canonical deterministic source and noise assumptions;

    \item the controlled randomized-regression policy;

    \item the boundary between logical signal-frame semantics and future AXI
    transport semantics.
\end{enumerate}

Evidence provenance is recorded in the SP-G0 detailed gate record and its
companion canonical-convention decision record.

\subsection{Acceptance Criteria}

SP-G0 may be marked PASS only when all of the following are satisfied:

\begin{enumerate}
    \item \(N\), \(M\), and logical \(\mathbf X\) indexing are frozen;

    \item covariance normalization and complex conjugation are frozen;

    \item mean treatment is explicit;

    \item ULA spacing and normalized sensor coordinates are explicit;

    \item angular reference and positive-angle direction are explicit;

    \item the steering trigonometric convention is explicit;

    \item the steering exponential sign is explicit;

    \item deterministic reference-case semantics are explicit;

    \item randomized regression policy is explicit;

    \item the SP-G0 versus INT-G0 responsibility boundary is explicit;

    \item historical convention conflicts are documented rather than hidden;

    \item SP-G1 can generate the canonical test case without making an
    undocumented convention choice.
\end{enumerate}

The completed SP-G0 pre-closure audit evaluated all twelve criteria and
confirmed that each acceptance criterion passed.

\subsection{Engineering Interpretation}

SP-G0 is deliberately treated as an engineering gate rather than as a
documentation exercise.

An array-processing implementation can produce numerically plausible results
while still being physically inconsistent because of a sign convention,
angle-reference convention, indexing convention, conjugation error, or frame
ordering mismatch.

Freezing these definitions before numerical optimization provides a controlled
reference against which later fixed-point and FPGA discrepancies can be
localized.

\subsection{Gate Status}

\textbf{Status: PASS.}

SP-G0 is closed.

The source audit and explicit project-decision process resolved the geometry,
steering, deterministic-case, mean-treatment, logical-frame, and randomized
regression conventions required by the project.

All twelve SP-G0 acceptance criteria passed. Repository diff hygiene, report
alignment, and the complete LaTeX smoke build also passed.

The authoritative Gate Registry and continuation handoff have been
synchronized with this closure.

SP-G1 remains PENDING and is the next permitted signal-processing gate after
this SP-G0 closure is integrated into the main branch.


\section{SP-G1 --- Array Signal and Snapshot Model}
\textbf{Purpose.} Establish the physical and mathematical meaning of the $8\times128$
complex input frame.\\
\textbf{Gate question.} What does one snapshot represent, how does the ULA steering law
produce inter-sensor phase structure, and can the deterministic frame be regenerated from
frozen conventions?\\
\textbf{Required theory.} Narrowband array model $\vec x[m]=\mat A(\boldsymbol\theta)\vec s[m]+\vec n[m]$,
steering vector, sensor coordinates, snapshot matrix.\\
\textbf{Evidence.} Deterministic single-source steering-vector case, dimensional checks,
phase-progression check.\\
\textbf{Exit criterion.} The canonical input frame is regenerated from documented
parameters and conventions.\\
\textbf{Status.} PENDING.

\section{SP-G2 --- Floating-Point SCM Golden Reference}
\textbf{Purpose.} Establish the authoritative software covariance result.\\
\textbf{Gate question.} Does an independent floating-point reference reproduce the same
SCM for common vectors and preserve expected matrix structure?\\
\textbf{Required theory.}
$\widehat{\mat R}=M^{-1}\mat X\mat X^{\mathsf H}$, element-wise conjugate products,
Hermitian structure, real-valued diagonal in exact arithmetic, positive-semidefinite behavior.\\
\textbf{Evidence.} Saved $\mat X$, saved $\widehat{\mat R}_{fp}$, independent NumPy/reference
comparison, Hermitian residual, diagonal checks.\\
\textbf{Exit criterion.} One deterministic and a small randomized set reproduce within
defined floating-point tolerance.\\
\textbf{Status.} PENDING.

\section{SP-G3 --- Statistical versus Numerical Error Separation}
\textbf{Purpose.} Prevent finite-snapshot estimation error from being misidentified as
fixed-point error.\\
\textbf{Gate question.} Which discrepancy belongs to statistical estimation and which to
finite-word-length representation?\\
\textbf{Required theory.}
\begin{equation}
\mat R_x \rightarrow \widehat{\mat R}_{fp} \rightarrow \widehat{\mat R}_{fxp},
\end{equation}
with $\widehat{\mat R}_{fp}-\mat R_x$ treated separately from
$\widehat{\mat R}_{fxp}-\widehat{\mat R}_{fp}$.\\
\textbf{Evidence.} Defined notation, one controlled finite-snapshot illustration, and code
variables that keep the two errors distinct.\\
\textbf{Exit criterion.} Statistical and numerical error metrics cannot be silently mixed.\\
\textbf{Status.} PENDING.

\section{SP-G4 --- Numeric Range and Fixed-Point Requirements}
\textbf{Purpose.} Convert signal properties into defensible word-length requirements.\\
\textbf{Gate question.} What input, product, and accumulator ranges are required under the
approved frame envelope?\\
\textbf{Required theory.} Quantization step, signed fixed point, scaling, complex
conjugate product, word growth, rounding, saturation, and overflow. For $M=128$, the
accumulation depth must be handled explicitly.\\
\textbf{Evidence.} Observed input envelope, two candidate formats, product/accumulator
width derivation, quantization policy, overflow/saturation instrumentation.\\
\textbf{Exit criterion.} Two candidates are fully specified without relying on implicit
compiler width behavior.\\
\textbf{Status.} PENDING.

\section{SP-G5 --- Fixed-Point Covariance Fidelity}
\textbf{Purpose.} Quantify how much the candidate numerical representation perturbs SCM.\\
\textbf{Primary metric.}
\begin{equation}
\epsilon_R=\frac{\|\widehat{\mat R}_{fxp}-\widehat{\mat R}_{fp}\|_F}
{\|\widehat{\mat R}_{fp}\|_F}.
\end{equation}
Additional evidence includes maximum element error, Hermitian residual, diagonal
consistency, and overflow/saturation counts.\\
\textbf{Exit criterion.} Each candidate has a traceable fidelity record on deterministic and
regression vectors.\\
\textbf{Status.} PENDING.

\section{SP-G6 --- Downstream DOA Consistency}
\textbf{Purpose.} Convert covariance error into an application-level consequence without
moving eigendecomposition or MUSIC into FPGA logic.\\
\textbf{Gate question.} Does the fixed-point covariance materially alter the existing host
DOA estimate for the approved test set?\\
\textbf{Evidence.}
$\widehat{\mat R}_{fp}\rightarrow\widehat\theta_{fp}$ and
$\widehat{\mat R}_{fxp}\rightarrow\widehat\theta_{fxp}$ with
$\Delta\theta=\widehat\theta_{fxp}-\widehat\theta_{fp}$. Eigenvalue/eigenspace diagnostics
are added only if needed to explain a discrepancy.\\
\textbf{Exit criterion.} Application-level numerical effect is quantified and reproducible.\\
\textbf{Status.} PENDING.

\section{SP-G7 --- Signal-Processing Acceptance}
\textbf{Purpose.} Close the numerical design from a signal-processing perspective before
hardware cost is used to select the final implementation.\\
\textbf{Decision basis.} Covariance fidelity, overflow behavior, structural preservation,
and downstream DOA consistency.\\
\textbf{Exit criterion.} A documented DSP recommendation is issued for the format to be
handed to FPGA-G0; any unresolved limitation is explicit.\\
\textbf{Status.} PENDING.
